\section{模型即产品，数据即模型：Agentic Web 与投资主线}

上一章说 Online Learning 需要真实交互和反馈，本章把这个判断落到产品、数据和投资策略。广密反复强调“模型即产品，数据即模型”。这句话看似绕口，其实非常具体：过去几年最大的体验升级来自模型升级，而模型升级的根部又来自数据分布升级。谁能拿到差异化数据，谁就能在模型和产品上形成差异。

\begin{figure}[H]
\centering
\includegraphics[width=0.96\textwidth]{model-data-product-loop.png}
\caption{模型即产品，数据即模型：用户任务、交互数据、专家蒸馏、模型升级和产品体验构成闭环。自制概念图，依据 00:43:05--00:56:45 对谈内容整理。}
\end{figure}

\begin{knowledgebox}{读图：产品体验回到数据分布}
用户任务暴露真实工作流，交互数据记录偏好和失败点，专家蒸馏补足高质量过程，模型升级扩大能力边界，产品体验改善后又吸引更多任务和反馈。这个闭环就是“模型即产品，数据即模型”的实际含义。
\end{knowledgebox}

\subsection{一横一纵：专家知识蒸馏与 Online Learning}

广密把未来路径概括为“一横一纵”。横向，是蒸馏人类专家知识，把更多行业、任务和工具链纳入模型的数据分布；纵向，是 Online Learning、Memory、主动 Agent 和新交互形态。横向可以让 AI 覆盖更多白领工作，纵向则可能让 AI 从用户和环境中持续成长。

\begin{importantbox}{一横一纵}
横向解决“还不会哪些行业和任务”，靠专家数据、真实工作流和行业场景扩展能力边界；纵向解决“能否越用越聪明”，靠 Memory、Online Learning、主动 Agent 和长期交互。
\end{importantbox}

这对创业公司有直接启发。只做一个薄 UI，很容易被基础模型公司或大平台复制；但如果能获得独特数据、专家过程、用户反馈和行业工作流，就可能在基础模型之上形成“内容资产”或“过程资产”。这也是为什么本期会反复讨论 Coding、Office、Finance、医疗、材料、机器人等具体领域。

\subsection{Agentic Web：真正的 Web3？}

节目中另一个重要概念是 Agentic Web。这里的 Web3 不是区块链意义上的 Web3，而是指 Agent 成为新入口后，网页、平台、支付、搜索和用户任务之间的关系被重组。过去互联网的核心入口是搜索框、App、Feed 和 Super App；Agentic Web 的核心入口可能是“我想完成一件事”。

如果 Agent 可以持续运行，它就不仅是内容消费工具，而是分配注意力、订单、流量和交易的中介。对 Uber、携程、搜索广告、手机厂商和超级应用来说，这都是重新谈判入口权的变量。

\begin{warningbox}{Agentic Web 的基础设施缺口}
Agentic Web 不是只靠模型能力就能落地。它还需要身份、授权、支付、风控、人机验证、网站协议、审计日志和责任归属。如果这些配套跟不上，Agent 会被卡在“能演示，难规模化执行”。
\end{warningbox}

\subsection{投资策略：押技术成长最陡峭的地方}

前面已经把产品、数据和 Agentic Web 连成闭环，本节转向投资框架。这里讨论投资不是为了给出买卖建议，而是用投资人的分类方式反向理解产业结构：钱会流向哪一层，通常意味着哪一层最稀缺、最能放大技术变化。

广密把 AI 投资策略概括为：投最领先模型、投最领先模型需要的算力和硅基基础设施、投最领先模型技术溢出的红利。这个框架对应三类资产：OpenAI/Google/Anthropic/字节这样的模型与入口公司，Nvidia/TSMC/算力基础设施这样的底层供给，以及 Cursor、Perplexity、Coding Agent、数据标注、Serving、多模态工具等技术溢出型公司。这里的 Perplexity 指 AI 搜索/问答公司名，不是语言模型评测中的 perplexity/PPL；后者是交叉熵指数化后的“困惑度”指标，用来粗略表示模型面对下一个 token 时的平均不确定性。

节目中还提到一个组合变化：从更集中押 OpenAI 与字节，转向 OpenAI、字节、Google、Anthropic、Nvidia、TSMC 等构成的篮子。背后的逻辑是交替领先和范式不确定性：如果很难判断最终赢家，就应该围绕主线做组合，而不是只押单一公司。

\begin{knowledgebox}{这不是投资建议，而是产业地图}
本讲义保留节目中的投资框架，是为了理解产业结构：模型、算力、平台入口、应用溢出和数据服务如何相互拉动。它不是对任何证券或公司估值的推荐。
\end{knowledgebox}

\subsection{2026 的主题：Online Learning、知识工作者、多模态与 Agentic Web}

本节把前面的技术和产业判断收束到 2026 年的观察清单。与其问“谁会最终赢”，更可操作的问题是：哪些主题会先出现可验证信号，哪些场景会先产生收入和用户行为变化。

面向 2026，节目列出几个期待主题。第一是 Online Learning 或下一范式信号；第二是更主动的 Agent 和更强 Memory；第三是围绕知识工作者的局部 L3/L4 能力，例如 Coding、Office/PPT/Excel、金融投研和长尾信息获取；第四是多模态大年；第五是 Agentic Web 对流量分配权的冲击。

这里的 L3/L4 是借用自动驾驶分级的类比：L3 可以理解为在特定条件下系统接管主要任务，人类仍需监督或接管；L4 是在限定场景内高度自动完成。广密认为“局部 L3/L4”很现实，“整体 L4”仍然很难。换句话说，先在清晰工作流中吃到经济价值，而不是等待全能 AGI。

\subsection{本章小结}

“模型即产品，数据即模型”把技术、产品和投资连成一个闭环。模型体验来自数据，数据来自真实任务，真实任务来自产品入口，产品入口又带来更多反馈。2026 的关键不是谁说自己更通用，而是谁能在具体场景里形成数据和产品的复利。

